\documentclass[11pt,a4paper]{article}
\usepackage{amsmath,amssymb}
\usepackage{enumitem}
\usepackage{booktabs}
\usepackage{geometry}
\usepackage{hyperref}
\usepackage{microtype}
\geometry{margin=2.5cm}
\setlength{\parskip}{6pt}
\setlength{\parindent}{0pt}

% ============================================================
% DOMAIN AUDIT
% ============================================================
% DOMAIN PROFILE USED: topology_optimization_eigenfrequency v1.0
%
% DOMAIN-SPECIFIC CHECKS COMPLETED:
%   spurious_low_density_modes, simp_mass_interpolation,
%   grey_regions_in_topology, frozen_eigenvector_approximation,
%   repeated_eigenvalues_mode_coalescence, mode_switching,
%   design_dependent_load_sensitivity, benchmark_implementation_disclosure
%
% DOMAIN-SPECIFIC CHECKS NOT POSSIBLE:
%   Visual inspection of topology figure grey-region extent (text-based
%   analysis only; density arrays not available to reviewer).
%   Verification of spurious low-density modes during iteration
%   (source code and intermediate eigenspectrum not available).
%
% DOMAIN_UNVERIFIED WARNINGS:
%   none (DOMAIN_UNVERIFIED = false, match_confidence = HIGH)
% ============================================================

\begin{document}

% ============================================================
% HEADER
% ============================================================

\begin{center}
{\Large\bfseries Peer Review Report}\\[6pt]
{\large Springer Nature -- Engineering with Computers}\\[4pt]
{\normalsize Manuscript: \textit{Frequency maximization of planar structures
using quasi-static approximation in topology optimization}}\\[4pt]
{\normalsize Review produced by: Multimodel AI peer review pipeline
(GPT\,+\,Claude, consensus synthesis)\\
Date of review: 2026-05-28}
\end{center}

\bigskip
\hrule
\bigskip

% ============================================================
% RECOMMENDATION
% ============================================================

\section*{Recommendation}

\textbf{MAJOR REVISION} \quad (Confidence: HIGH)

Both independent review branches (GPT and Claude) reached this decision independently
with high confidence and no contradictions on the core findings. The manuscript describes
a practically useful computational method, but four critical-level issues and numerous
major issues must be resolved before publication.

% ============================================================
% SUMMARY
% ============================================================

\section{Summary}

This paper proposes a quasi-static approximation for natural-frequency maximization in
SIMP topology optimization. The fundamental eigenpair is computed once on the initial
uniform-density domain and then kept fixed throughout the optimization, so that each
subsequent iteration reduces to a standard compliance minimization step. The load-
sensitivity term arising from the density-dependent mass matrix is omitted. A scalar
load-weighting parameter combines two modal inertial loads for simultaneous two-mode
frequency design. The method is benchmarked on a simply supported beam (comparison
against Olhoff and Yuksel implementations), a clamped beam (multi-mode demonstration),
and a building-like structure with passive elements.

The computational efficiency claim is verified: the reported speedups
($7.1\times$ over the Olhoff implementation and $2.7\times$ over the Yuksel method at
$400\times50$ mesh) are arithmetically consistent with Table~1, and the frequency gap of
$8.6\%$ below the Olhoff optimum is correctly computed. The paper has a publishable core.
However, the current version contains four critical methodological failures that prevent
acceptance.

% ============================================================
% CRITICAL ISSUES
% ============================================================

\section{Critical Issues}

\subsection*{C1.\; Unvalidated Load-Sensitivity Omission (Eq.~6)}

The inertial load $f(\mathbf{x}) = \omega_1^{(0)2} \mathbf{M}(\mathbf{x})\boldsymbol{\Phi}_1$
(Eq.~5) depends on element densities through the evolving mass matrix~$\mathbf{M}(\mathbf{x})$.
The total derivative of compliance with respect to~$x_e$ therefore includes a load-sensitivity
term
\[
  \frac{\partial \mathbf{f}}{\partial x_e}
    = \bigl(\omega_1^{(0)}\bigr)^2 \frac{\partial \mathbf{M}}{\partial x_e}\boldsymbol{\Phi}_1
    = \bigl(\omega_1^{(0)}\bigr)^2\, d\, x_e^{d-1}\, \mathbf{m}_e^0\, \boldsymbol{\Phi}_1.
\]
Eq.~(6) omits this term entirely, retaining only the stiffness-based sensitivity.
The justification---citing Munro~(2024) by analogy to self-weight compliance---is
heuristic and does not constitute an error bound or derivation specific to the
eigenfrequency case. No finite-difference gradient check is provided.

\textbf{Required revision:} Provide a finite-difference validation of Eq.~(6) against
the full gradient for at least one example. Alternatively, derive an analytical bound
showing that the omitted term is negligible relative to the retained stiffness sensitivity
across the relevant density range~$x_e\in[x_{\min},1]$.

\subsection*{C2.\; Inconsistent Frozen-Eigenpair Reference Design}

Three different reference designs are specified for the frozen eigenpair:
\begin{itemize}[nosep]
  \item Section~3.2: ``eigenvalue problem on the initial uniform-density field $x_e = V_f$
    for all~$e$''.
  \item Section~3.3: ``the first two eigenpairs of the initial solid domain''.
  \item Section~4.3: ``the first two eigenvectors of the fully solid design domain''.
\end{itemize}
For the building example with $V_f=0.3$, the uniform-density design
($x_e = 0.3$ everywhere) and the fully solid domain ($x_e = 1$) produce substantially
different stiffness and mass matrices and different eigenpairs. The reference
eigenfrequencies $\omega_1^{(0)} = 19.84$~rad/s and $\omega_2^{(0)} = 90.93$~rad/s are
stated for the building (Section~4.3) but the density at which they were computed is not
specified. The method is underspecified.

\textbf{Required revision:} Define the reference design precisely (density value and
equation) and apply it consistently across all sections and all examples. Report the
reference eigenfrequency explicitly for each example alongside the density used.

\subsection*{C3.\; Algorithm--Implementation Inconsistency and Contribution~(i)}

Algorithm~3.5 (step~e) specifies an ``Optimality Criteria (OC) bisection update.''
Sections~4.2 and~4.3 explicitly state that MMA is used with move limit~0.2. OC bisection
and MMA are different algorithms with different software dependencies and per-iteration
cost structures.

The abstract and contribution~(i) claim the method ``integrates directly into a standard
compliance minimization code without requiring an additional optimizer.'' MMA is a
non-trivial mathematical programming solver that is not part of a standard compliance
code. This claim holds only for the simply supported beam (OC); it does not hold for the
multi-mode examples.

\textbf{Required revision:} Either correct Algorithm~3.5 to reflect the actual optimizer
used per example (OC for the single-mode case, MMA for multi-mode cases), or restrict
contribution~(i) to the single-mode OC formulation and explicitly acknowledge that the
multi-mode examples require MMA.

\subsection*{C4.\; Non-Monotone Data in Table~3 Contradicts Contribution~(iii)}

Contribution~(iii) states: ``MAC-based demonstration that load-weighted modal gains scale
monotonically with the assigned mode weight.'' Section~4.2 (lines~599--608) asserts that
the $\widetilde{\omega}_2/\omega_2^{(0)}$ gain ``increases monotonically as $\alpha$
decreases.''

Table~3 reports the following $\Phi_2$-tracked gains as $\alpha$ decreases from~1.0
to~0.0:
\[
  1.99\times,\quad 1.73\times,\quad 2.05\times,\quad 2.25\times,\quad 2.45\times.
\]
The step from $\alpha=1.0$ to $\alpha=0.75$ is a \textbf{13\% decrease}
($1.99\to1.73$), not an increase. This is confirmed by arithmetic from the
tabulated frequencies: $629.1/362.8 = 1.73\times$ and $722.0/362.8 = 1.99\times$
(with $\omega_2^{(0)}\approx362.8$~rad/s implied by back-calculation from all five
table rows). The individual values in Table~3 are internally consistent; the monotone
claim is incorrect.

In addition to this arithmetic failure, even in cases where the MAC-tracked gain is
monotone, the target mode frequently appears at positions~33, 46, 71, 81, or~91 in the
final spectrum while lower modes remain below it. This demonstrates mode-shape-specific
stiffening, not simultaneous maximization of the fundamental frequency spectrum.

The building example (Table~5) does exhibit monotone $\Phi_1$ and $\Phi_2$ gain
sequences, and monotone $\Phi_1$ gain is verified for the clamped beam.

\textbf{Required revision:} Either (a)~correct the Table~3 data if a computation error
occurred, or (b)~restrict contribution~(iii) to the cases where monotonicity holds and
provide a physical explanation for the non-monotone step at $\alpha=0.75$ in the
clamped-beam $\Phi_2$ sequence. In either case, reframe the multimode contribution as
``MAC-tracked modal-gain stiffening'' rather than ``simultaneous multi-mode frequency
maximization,'' and report the full eigenspectrum below the tracked target mode for
representative $\alpha$ values.

% ============================================================
% MAJOR ISSUES
% ============================================================

\section{Major Issues}

\subsection*{M1.\; Speedup Comparison Against a Modified Olhoff Implementation}

Section~5.1 discloses that the authors' Olhoff implementation departs from
Olhoff and~Du~(2014) in three material respects:
(i)~a seven-level Heaviside projection continuation schedule,
(ii)~a grayness penalty term in the MMA objective, and
(iii)~a trial eigensolve after every MMA update, which doubles sparse Cholesky
factorizations per outer iteration.
These additions substantially inflate per-iteration cost relative to the canonical bound
formulation. The abstract and conclusions state ``$7.1\times$ faster than the bound
formulation'' without qualification.

\textbf{Required revision:} Label the speedup explicitly as being over the \emph{authors'
modified} Olhoff implementation, not over the canonical algorithm of Olhoff and
Du~(2014). Provide an estimate of the speedup that would be achieved over the original
formulation (without the extra trial eigensolve and projection chain), for example by an
ablation run or by analytical estimation of the removed overhead.

\subsection*{M2.\; Frequency Convergence with Mesh Refinement Not Reported}

Table~1 provides timing, iteration count, and RAM for four mesh resolutions but does not
report the achieved fundamental frequency~$\omega_1$ for each mesh and each method.
Without frequency-vs-mesh data it is not possible to assess whether the proposed method
converges to a mesh-independent frequency optimum.

\textbf{Required revision:} Add the achieved $\omega_1$ for each method and each mesh
resolution to Table~1 and discuss the convergence trend.

\subsection*{M3.\; No Reference Comparison for Multi-Mode Examples}

The clamped-beam and building examples (Sections~4.2--4.3) present no comparison against
the Olhoff or Yuksel method. Frequency gains are reported only relative to the initial
(unoptimized) solid domain, which is a trivial baseline.

\textbf{Required revision:} Provide at least one reference comparison (Olhoff or Yuksel)
for the clamped-beam or building example at a representative~$\alpha$ value, or justify
why such a comparison is not feasible.

\subsection*{M4.\; Missing Hardware Specification and Timing Variance}

Table~1 timing data are stated as averaged over 10~independent runs, but no standard
deviations are reported. CPU model, clock speed, RAM, operating system, MATLAB~2025b
build number, and BLAS/thread configuration are absent.

\textbf{Required revision:} Report standard deviations (or min/max) for Table~1 timing
entries. Add a hardware and software specification table.

\subsection*{M5.\; Emerged Low-Frequency Modes Not Characterized}

At $\alpha=1.0$ in the clamped-beam example, topo-modes~1 and~2
($\omega_1=275.28$~rad/s, $\omega_2=365.78$~rad/s from Table~2) exhibit negligible MAC
correlation with any solid-domain mode (maximum MAC$<0.04$). The paper notes this but
provides no mode-shape plots for these modes. Without the mode shapes it cannot be
determined whether they are physically meaningful structural modes or artifacts.

\textbf{Required revision:} Provide mode-shape plots for topo-modes~1 and~2 at
$\alpha=1.0$ for the clamped-beam example and discuss their physical interpretation.

\subsection*{M6.\; Discreteness Metrics Absent}

The Olhoff implementation enforces grayness${}<0.05$ as a convergence criterion, but the
proposed method applies no Heaviside projection and no grayness criterion. No
grey-element fraction or equivalent discreteness metric is reported for any
proposed-method topology.

\textbf{Required revision:} Report a discreteness metric
(e.g.,~$g = n_e^{-1}\sum_e 4x_e(1-x_e)$) for all proposed-method topologies in the
figure captions or a supplementary table.

\subsection*{M7.\; Spurious Low-Density Modes Claim Unsubstantiated}

Section~3.1 states that the proposed method uses $d=p=3$ ``without inducing localized
modes in the examples reported here.'' The Olhoff implementation uses $d=6$ for
$x_e\leq0.1$ specifically to suppress spurious low-frequency modes
(Du and Olhoff~2007). The assertion that $d=3$ is safe is not supported by any
eigenspectrum evidence.

\textbf{Required revision:} Provide evidence that spurious low-density modes do not arise
(e.g.,~eigenspectrum inspection of converged topologies or minimum-frequency monitoring
during optimization).

\subsection*{M8.\; Introduction Gap Claim Inadequately Bounded Against Huang et~al.~(2025)}

The introduction (lines~101--103) states that the multi-mode extension ``has not been
systematically demonstrated.'' However, Huang et~al.~(2025), cited in Section~2.3,
presents a MAC-constrained multi-mode topology optimization approach with mode-tracking
constraints during iteration.

\textbf{Required revision:} Revise the gap claim to acknowledge Huang et~al.~(2025) and
explain specifically how the proposed scalar-weighted inertial load differs from their
MAC-constrained, within-iteration multi-mode framework.

% ============================================================
% EQUATION AND MATHEMATICAL REVIEW
% ============================================================

\section{Mathematical and Equation Review}

\begin{enumerate}[label=\textbf{E\arabic*.},nosep]

\item \textbf{Eq.~(6) gradient omission.} As stated under C1, the omitted term
  $(\omega_1^{(0)})^2 d\, x_e^{d-1} \mathbf{m}_e^0 \boldsymbol{\Phi}_1$ is
  $O((\omega_1^{(0)})^2 d\, x_e^{d-1})$, which need not be small relative to
  the retained stiffness sensitivity at intermediate densities. No bound is given.

\item \textbf{Eq.~(7) load-scale asymmetry.}
  The combined load $f_\alpha = \alpha(\omega_1^{(0)})^2\mathbf{M}\boldsymbol{\Phi}_1
  + (1-\alpha)(\omega_2^{(0)})^2\mathbf{M}\boldsymbol{\Phi}_2$ scales the two
  contributions by $(\omega_j^{(0)})^2$. For the clamped beam, the ratio
  $(\omega_2^{(0)}/\omega_1^{(0)})^2\approx6.2$, so $\alpha=0.5$ is not mechanically
  symmetric between the two modes. This scaling choice is not discussed.

\item \textbf{Eq.~(9) MAC formula.} The formula uses Euclidean inner products rather
  than mass-weighted inner products. Comparing eigenvectors from different density
  fields (initial solid vs.\ final topology) with Euclidean MAC is non-standard; the
  mass-weighted version is more physically meaningful. The choice should be stated.

\item \textbf{Algorithm~3.5, step~2.} The parameter $n_\text{modes}$ is not defined.
  Specify $n_\text{modes}=1$ for the simply supported beam and $n_\text{modes}=2$ for
  the multi-mode examples.

\item \textbf{Dimensional consistency.} Eq.~(5) through Eq.~(7): units of
  $f(\mathbf{x})$ are (rad/s)$^2 \times$ kg/m $\times$ (dimensionless) = N/m (2D
  unit-thickness). Consistent with compliance $c = \mathbf{u}^\top\mathbf{K}\mathbf{u}$
  having units N$\cdot$m. Verified~\checkmark.

\end{enumerate}

% ============================================================
% REPRODUCIBILITY
% ============================================================

\section{Reproducibility Concerns}

\begin{enumerate}[label=\textbf{R\arabic*.},nosep]

\item \textbf{Hardware and software.} CPU model, clock speed, RAM, OS version,
  MATLAB 2025b patch level, required toolboxes, and BLAS thread count are absent.

\item \textbf{Timing variance.} Ten runs averaged in Table~1; no standard deviations
  or source of variance stated. For OC bisection with fixed parameters the algorithm
  is deterministic; the measurement methodology should be described.

\item \textbf{Repository.} No release tag or commit hash cited. No license stated.
  No per-table or per-figure script mapping. Required MATLAB toolboxes not listed.

\item \textbf{Convergence tolerance inconsistency.} $\varepsilon_\mathrm{tol}=10^{-3}$
  for simply supported beam (Algorithm~3.5); $\varepsilon_\mathrm{tol}=2\times10^{-3}$
  for the building example (Section~4.3). Justify the difference or use a uniform
  tolerance.

\item \textbf{Void pseudo-material inconsistency.} $\rho_\mathrm{min}=10^{-9}\rho_0$
  for the simply supported beam (Section~4.1); $\rho_\mathrm{min}=10^{-6}\rho_0$ for
  the clamped beam and building (Sections~4.2--4.3). The $1000\times$ difference is
  not discussed.

\end{enumerate}

% ============================================================
% REQUIRED REVISIONS
% ============================================================

\section{Required Revisions Summary}

\begin{enumerate}[label=\textbf{\arabic*.},nosep]
  \item Provide finite-difference gradient validation for Eq.~(6) or an analytical
    error bound on the omitted load-sensitivity term.
  \item Unify the frozen-eigenpair reference design to one precisely stated density
    and apply consistently across all sections and examples.
  \item Reconcile Algorithm~3.5 and Sections~4.2--4.3 regarding the update rule
    (OC vs.\ MMA) and revise contribution~(i) accordingly.
  \item Correct the Table~3 $\Phi_2$-gain monotone claim or correct the data; provide
    an explanation for the non-monotone step at $\alpha=0.75$.
  \item Reframe contribution~(iii) and the multi-mode results as MAC-tracked
    shape-stiffening; report the full eigenspectrum below the tracked target mode.
  \item Label the speedup as being over the \emph{authors' modified} Olhoff
    implementation; estimate or bound the speedup over the canonical algorithm.
  \item Add the achieved $\omega_1$ for each mesh resolution to Table~1 (frequency
    vs.\ mesh convergence).
  \item Add at least one reference comparison for the clamped-beam or building example.
  \item Add hardware specification and timing standard deviations.
  \item Provide mode-shape plots for topo-modes~1 and~2 at $\alpha=1.0$
    (clamped-beam example).
  \item Report a discreteness metric for all proposed-method topologies.
  \item Demonstrate absence of spurious low-density modes for $d=p=3$.
  \item Revise the introduction gap claim regarding Huang et~al.~(2025).
\end{enumerate}

% ============================================================
% MINOR COMMENTS
% ============================================================

\section{Minor Comments}

\begin{enumerate}[label=\textbf{m\arabic*.},nosep]
  \item Section~4.1 refers to ``Figs.~5--7'' for the simply supported beam topologies;
    these are Figs.~1--3. Correct the reference.
  \item SIMP interpolation (Eq.~1) lacks its canonical citation
    (Bends{\o}e and Sigmund~1999 or equivalent). Add at first use.
  \item The Heaviside projection filter used in the Olhoff implementation
    (Section~5.1) lacks a citation (Wang et~al.~2011 or equivalent).
  \item Rayleigh's principle is invoked as a theoretical foundation without a canonical
    structural dynamics reference (e.g.,~Courant and Hilbert~1953).
  \item Report the clamped-beam solid-domain reference eigenfrequencies
    $\omega_1^{(0)}\approx145.4$~rad/s and $\omega_2^{(0)}\approx362.8$~rad/s directly
    (currently back-calculable from Table~3 but not stated).
  \item Define $n_\mathrm{modes}$ explicitly in Algorithm~3.5.
  \item The Yuksel and Yilmaz~(2025) reference is published in
    \textit{Engineering Computations}; confirm the citation details against the journal
    record.
  \item The Li et~al.~(2021) reference lacks a DOI; add for consistency.
  \item Clarify the non-standard use of ``quasi-static'' (standardly: low-rate dynamic
    limit; here: static compliance under frozen inertial load) to avoid
    misinterpretation.
\end{enumerate}

% ============================================================
% META-REVIEW NOTE
% ============================================================

\section*{Cross-Review Note}

This report synthesizes independent reviews from two AI review pipelines (GPT and Claude).
All critical and major findings are supported by direct evidence from the manuscript
(equation numbers, section numbers, or table entries).
The arithmetic in Sections~4.1, 4.3 (speedups, frequency gap, building gains) was
independently verified and found correct. The non-monotone $\Phi_2$-gain sequence in
Table~3 was confirmed arithmetically from the tabulated frequencies.
No hallucinated claims were identified in either review branch.
Discarded signals: absence of scale bars (not standard in topology optimization); RAMP
citation (RAMP not used in this paper).

% ============================================================
% CONSENSUS GAPS
% ============================================================

\section*{Known Coverage Gaps}

The following issues were identified as partially outside the scope of automated review
and should receive domain-expert attention during revision assessment:

\begin{enumerate}[label=\textbf{G\arabic*.},nosep]
  \item \textbf{Initialization eigensolve timing.} The wall time of the single
    initialization eigensolve (step~2 of Algorithm~3.5) is not reported separately
    from the per-iteration compliance solve. For design scenarios with repeated
    re-solves from the same initial domain, this cost should be disclosed.
  \item \textbf{Frozen eigenvector accuracy near mode coalescence.} For the simply
    supported beam benchmark the Olhoff optimum is bimodal ($\omega_1$ and~$\omega_2$
    coalesce). The behavior of the frozen-eigenvector approximation as the two
    frequencies approach each other during optimization is unanalyzed. This is
    particularly relevant because the frozen mode from the initial uniform design may
    become increasingly inaccurate as the topology approaches the bimodal state.
\end{enumerate}

\end{document}
